
\usepackage{graphicx}
\usepackage{latexsym}
\usepackage{listings}
\usepackage{amsmath}
\usepackage{multirow}
\usepackage{multicol}
\usepackage{amsfonts}
\usepackage{amssymb}
\usepackage{epsfig,psfrag} 
\usepackage{wrapfig}
\usepackage[left=1.8cm,top=1cm,right=1.8cm,bottom=1cm]{geometry}
\usepackage[lined, ruled, vlined]{algorithm2e}
\usepackage[normalem]{ulem}
\usepackage{graphics}
\usepackage{pstricks,pst-node,pst-tree}
\usepackage{eqnarray,amsmath}
\usepackage{caption}
\usepackage{circuitikz}
\usepackage{datetime}
\usepackage{nicematrix, tikz}
\usepackage{mathtools}
\ctikzset{bipoles/thickness=1}
\usepackage[colorlinks=true, allcolors=blue]{hyperref}
\usepackage{listings}
\usepackage{xcolor}
\usepackage{ragged2e}
\usepackage{enumitem}
\usepackage{ifthen}
\usepackage{minted}
\setminted{
    fontsize=\footnotesize,
    breaklines=true,
    tabsize=4
}
\usepackage{pythonhighlight}
\usepackage[utf8]{inputenc} 
\usepackage[shorthands=off,brazil]{babel}
\AtBeginDocument{\shorthandoff{"}}
\usepackage{ifthen}


\definecolor{cinza-codigo}{gray}{0.92}
\setlength{\fboxsep}{0.8pt}
\emergencystretch 2em

\ExplSyntaxOn
\cs_new_protected:Npn \__code_box:n #1 {
  \tl_if_empty:nF { #1 }
   {
     \colorbox{cinza-codigo}{\texttt{#1}}
   }
}

\NewDocumentCommand{\code}{m}
 {
   \seq_set_split:Nnn \l_tmpa_seq { ~ } { #1 }
   \int_set:Nn \l_tmpa_int { 0 }
   \seq_map_inline:Nn \l_tmpa_seq
    {
      \int_incr:N \l_tmpa_int
      \int_compare:nF { \l_tmpa_int == 1 } { \allowbreak\ \hspace{0pt} }
      
      \tl_set:Nn \l_tmpa_tl { ##1 }
      \tl_replace_all:Nnn \l_tmpa_tl { \_ } { \allowbreak \_ \allowbreak }
      \tl_replace_all:Nnn \l_tmpa_tl { ; } { ; \allowbreak }
      \tl_replace_all:Nnn \l_tmpa_tl { , } { , \allowbreak }
      \tl_replace_all:Nnn \l_tmpa_tl { : } { : \allowbreak }
      \tl_replace_all:Nnn \l_tmpa_tl { / } { / \allowbreak }
      \tl_replace_all:Nnn \l_tmpa_tl { . } { . \allowbreak }
      \tl_replace_all:Nnn \l_tmpa_tl { ( } { ( \allowbreak }
      \tl_replace_all:Nnn \l_tmpa_tl { ) } { ) \allowbreak }
      \tl_replace_all:Nnn \l_tmpa_tl { [ } { [ \allowbreak }
      \tl_replace_all:Nnn \l_tmpa_tl { ] } { ] \allowbreak }
      \tl_replace_all:Nnn \l_tmpa_tl { \{ } { \{ \allowbreak }
      \tl_replace_all:Nnn \l_tmpa_tl { \} } { \} \allowbreak }
      \tl_replace_all:Nnn \l_tmpa_tl { " } { \allowbreak " \allowbreak }
      \tl_replace_all:Nnn \l_tmpa_tl { = } { \allowbreak = \allowbreak }
      \tl_replace_all:Nnn \l_tmpa_tl { -> } { -> \allowbreak }
      \tl_replace_all:Nnn \l_tmpa_tl { < } { < \allowbreak }
      \tl_replace_all:Nnn \l_tmpa_tl { > } { > \allowbreak }
      \tl_replace_all:Nnn \l_tmpa_tl { + } { + \allowbreak }
      \tl_replace_all:Nnn \l_tmpa_tl { - } { - \allowbreak }
      \tl_replace_all:Nnn \l_tmpa_tl { * } { * \allowbreak }
      
      \seq_set_split:Nnn \l_tmpb_seq { \allowbreak } { \l_tmpa_tl }
      \int_set:Nn \l_tmpb_int { 0 }
      \seq_map_inline:Nn \l_tmpb_seq
       {
         \int_incr:N \l_tmpb_int
         \int_compare:nF { \l_tmpb_int == 1 } { \allowbreak }
         \__code_box:n { ####1 }
       }
    }
 }
\ExplSyntaxOff

% 2. Configurar o estilo do bloco de código
\lstset{
    language=bash,                  % Define a linguagem como Bash
    basicstyle=\ttfamily\small,     % Fonte monoespaçada e menor
    keywordstyle=\color{blue},      % Cor para comandos/palavras-chave
    commentstyle=\color{gray},      % Cor para comentários
    stringstyle=\color{purple},     % Cor para strings
    numbers=left,                   % Mostra números de linha na esquerda
    numberstyle=\tiny\color{gray},  % Estilo dos números de linha
    stepnumber=1,                   % Número em cada linha
    backgroundcolor=\color{lightgray!20}, % Cor de fundo do bloco
    frame=single,                   % Moldura ao redor do código
    breaklines=true,                % Quebra linhas longas automaticamente
    showstringspaces=false          % Não mostra espaços especiais em strings
}


% Estilização para blocos de código Python
% Configuração do estilo para o código Python
\lstdefinestyle{pythonstyle}{
    language=Python,
    basicstyle=\ttfamily\small,         % Fonte monoespaçada pequena
    keywordstyle=\color{blue}\bfseries, % Palavras-chave em azul e negrito
    stringstyle=\color{red},            % Strings em vermelho
    commentstyle=\color{green!50!black},% Comentários em verde escuro
    numbers=left,                       % Numeração das linhas no lado esquerdo
    numberstyle=\tiny\color{gray},      % Estilo dos números das linhas
    stepnumber=1,
    numbersep=10pt,
    backgroundcolor=\color{gray!10},    % Fundo cinza claro
    frame=single,                       % Moldura ao redor do código
    tabsize=4,
    showstringspaces=false,
    breaklines=true                     % Quebra linha automaticamente se for longa
}
\lstset{style=pythonstyle}



\newcommand\encircle[1]{%
  \tikz[baseline=(X.base)] 
    \node (X) [draw, shape=circle, inner sep=0] {\strut #1};}
\newdateformat{specialdate}{\THEYEAR-\twodigit{\THEMONTH}-\twodigit{\THEDAY}}
\extraheadheight[.25in]{.5in}
\extrafootheight[.25in]{.5in}
\pagestyle{empty}
\usepackage{xcolor}
\definecolor{greendel}{HTML}{2F5416}
\definecolor{bluepoli}{HTML}{1E1D44}
\definecolor{redpoli}{HTML}{C73834}

